% !TeX root = ../colloquium-2.tex

% Определение неопределённого интеграла и таблица интегралов.

\phantomsection\label{sec_Definition21}
\begin{definition}[
    Определение неопределённого интеграла и таблица интегралов.
]
    \subsection*{Определение неопределённого интеграла.}
    \begin{note}[purple]
        Произвольная первообразная функциии $f$, определённой на некотором интервале, называется 
        неопределённым интегралом функции $f$ и обозначается $\int f(x) \, dx = F(x) + C$. \\
        "Произвольная первообразная функции $f$ равна какой-то известной первообразной плюс 
        некоторая постоянная".
    \end{note}

    \subsection*{Таблица интегралов.}
    \begin{note}[purple]
        \begin{enumerate}
            \item 
        \end{enumerate}
    \end{note}
\end{definition}